\begin{abstract}
Efficient long-video understanding requires identifying relevant evidence while minimizing the tokens processed by vision-language models (VLMs). Existing video retrieval-augmented generation (RAG) methods rely on frame selection and token compression, but overlook how much each additional token is worth. We introduce \textbf{\name}, a plug-and-play framework that formulates video RAG as a token allocation problem over frames and subtitles under a shared budget. Our key insight is that the value of a token diminishes within each item and decreases further as similar evidence receives more tokens. \name captures both effects in a single objective and solves it with an efficient iterative water-filling algorithm that has a provable convergence guarantee. For multi-turn conversations, \name reuses the tokens processed in earlier turns and refines relevance with the VLM's attention at no extra cost. Our evaluations show that \name improves accuracy by up to 7.2 percentage points over state-of-the-art baselines at comparable compute. On multi-turn video understanding, \name achieves up to 1.5 percentage points higher accuracy with 1.7--17.3$\times$ fewer TFLOPs.
% Long video understanding with vision-language models (VLMs)
% requires balancing evidence coverage and detail in
% retrieval-augmented generation (RAG).
% Existing RAG approaches adapt evidence selection but do not
% explicitly optimize ``token value''---a token's estimated contribution
% to answer accuracy.
% We present \textbf{\name}, an efficient video RAG framework that allocates a
% shared token budget across evidence items (including frames and subtitles) to maximize aggregate token value.
% First, we model token value using query relevance,
% diminishing accuracy gains, and evidence redundancy.
% Second, we develop a token allocator that uses iterative water-filling
% to allocate tokens efficiently.
% Finally, we support multi-turn conversations through incremental allocation
% and attention refinement, reusing evidence and updating relevance as
% queries change.
% Across various long-video benchmarks and two VLM backbones, \name achieves the best accuracy-cost trade-off.
\end{abstract} 
